
\noindent\textbf{۱. نحوه بستن پتانسیومتر و منبع تغذیه DC چگونه است؟ چرا؟}
\newline
دو سر ثابت پتانسیومتر به‌صورت موازی با منبع $5\,\mathrm{V}$ بسته می‌شود و
خروجی از کنتاکت متحرک و یکی از دو سر ثابت گرفته می‌شود. این آرایش یک
تقسیم‌کننده‌ی ولتاژ است:
\(V_{\mathrm{out}} = V_{\mathrm{in}}\,\frac{R_x}{R}\)
\newline
که $R_x$ مقاومت بین کنتاکت متحرک و سر زمین‌شده است. با جابه‌جا کردن کنتاکت،
$R_x$ از $0$ تا $R$ تغییر می‌کند و ولتاژ پیوسته‌ی $0$ تا $5\,\mathrm{V}$
به دست می‌آید.
دلیل استفاده از این آرایش آن است که به ولتاژ معکوس پیوسته‌ی قابل تنظیم از
$0$ تا $5\,\mathrm{V}$ نیاز داریم؛ در حالی که با اتصال سری فقط جریان
تغییر می‌کند و ولتاژ خروجی هرگز به صفر نمی‌رسد.

\khat

\noindent\textbf{۲. آزمایش فوتوالکتریک مبتنی بر کدام خاصیت تابش الکترومغناطیسی است؟ توضیح دهید.}
\newline
این آزمایش بر ماهیت ذره‌ای نور استوار است.
تابش الکترومغناطیسی به صورت بسته‌های گسسته‌ی انرژی به نام فوتون منتشر و جذب می‌شود.
انرژی هر فوتون
\((E_{\gamma} = h\nu)\)
فقط به فرکانس آن وابسته است.
\newline
یک فوتون فقط با یک الکترون برهم‌کنش می‌کند و کل انرژی‌اش را به الکترون می‌دهد.
بخشی از این انرژی صرف غلبه بر تابع کار می‌شود و باقی آن، انرژی جنبشی الکترون خواهد بود.
\begin{equation*}
    K_\textrm{\lr{max}} = h \nu - \varphi
\end{equation*}
در آزمایش، ولتاژ معکوس آن‌قدر زیاد می‌شود که حتی پرانرژی ترین فوتو‌الکترون ها هم به آند نرسند، یعنی:
\[K_\textrm{\lr{max}} = eV \implies V = \frac{h}{e}\,\nu - \frac{\varphi}{e}\]
این همان رابطه‌ای هست که برازش خطی بر آن انجام شده.

\khat

\noindent\textbf{۳. مواردی که فرضیات کلاسیکی از عهده توجیه آن‌ها برنمی‌آیند برشمرده و شاهد آن را در این آزمایش بیان کنید.}
\begin{enumerate}[leftmargin=*, itemsep=0.8em]

    \phenomenon{وجود بسامد آستانه}
    {با شدت کافی، هر بسامدی باید بتواند الکترون آزاد کند.}
    {وقتی \(\nu < \nu_0\) است، هیچ فوتوالکترونی خارج نمی‌شود، هرچقدر هم که شدت زیاد باشد.}
    {خط برازش‌شده در \(\nu_0\) ولتاژ را صفر می‌کند و پایین‌تر از آن جریانی نداریم.}

    \phenomenon{وابستگی \(K_{\max}\) به بسامد، نه شدت}
    {شدت بیشتر یعنی میدان قوی‌تر، پس الکترون باید انرژی بیشتری بگیرد.}
    {\(K_{\max}\) فقط به بسامد بستگی دارد و شدت فقط تعداد الکترون‌ها را تغییر می‌دهد.}
    {هر رنگ ولتاژ متوقف‌کننده‌ی مخصوص خود را دارد و رابطه‌ی \(V-\nu\) خطی است.}

\end{enumerate}

\khat

\noindent\textbf{۴. آیا می‌توانید با توجه به مبانی فیزیک کوانتومی، پدیده‌هایی که رخ می‌دهد را توجیه کنید؟}
\newline
در فلز، الکترون‌ها در ترازهای انرژی مقید هستند و تابع کار \(\varphi\) فاصله‌ی
تراز فرمی تا تراز خلأ است. یک الکترون آزاد نمی‌تواند فوتون را کامل جذب کند،
زیرا انرژی و تکانه هم‌زمان پایسته نمی‌مانند.

در فلز، شبکه‌ی بلور تکانه‌ی
اضافی را می‌گیرد و جذب ممکن می‌شود. احتمال جذب تک‌فوتونی با شدت
متناسب است، اما جذب هم‌زمان دو فوتون با \(I^2\) متناسب است و در نور معمولی
ناچیز است؛ پس هر الکترون فقط از یک فوتون انرژی می‌گیرد.
فقط الکترون‌های نزدیک تراز فرمی به \(K_{\max}\) می‌رسند و بقیه انرژی کمتری
دارند؛ به همین دلیل جریان به‌تدریج به صفر می‌رسد و تعیین ولتاژ متوقف‌کننده
با خطا همراه است. همچنین رنگ‌های گسسته‌ی لامپ جیوه از گذار بین ترازهای
گسسته‌ی انرژی اتم جیوه می‌آیند.